\section{Proof Outline in Strongly Convex Linear Regression}\label{sec:convex-proof-outline}
In this section, we give the outline of \Cref{thm:convex-L-estimate}, \Cref{thm:convex-opt-lr}, and \Cref{thm:convex-E_K}. The main technical challenges and our proof insights are briefly stated in \Cref{sec:strongly-convex-tech}.

\Cref{sec:strongly_convex} centres on \Cref{thm:convex-E_K}, which establishes a scaling law for the effective reuse rate $E(K,N)$ in terms of the relative magnitudes of number of epochs $K$ and dataset size $N$. Its proof unfolds in three stages.

\paragraph{1. An explicit approximation of the expected excess risk.} \Cref{thm:convex-L-estimate} derives a sufficiently accurate asymptotic formula for the expected excess risk of multi-epoch SGD. The argument begins with a bias–variance decomposition, splitting the expected excess risk into a variance term (\Cref{lem:variance-estimate}) and a bias term (\Cref{lem:bias-estimate}).
\begin{itemize}
    \item \textbf{Variance term.} The closed-form approximation relies on concentration properties of matrix contractions together with a careful treatment of data shuffling.
    \item \textbf{Bias term.} The same contraction inequality is employed to obtain an analytic expression, after which tight error bounds are proved for the full range of relative sizes of $K$ and $N$. These bounds hold uniformly over a broad class of learning rates, necessitating detailed case-by-case analysis.
\end{itemize}

\paragraph{2. Selection of a nearly optimal learning rate.}
\Cref{thm:convex-opt-lr} identifies a learning rate whose resulting loss is asymptotically equivalent to the minimum excess risk attained with the optimal learning rate as stated in \Cref{sec:setup}. This ``approximately optimal learning rate'' will be fixed in Appendix~\ref{proof:convex-opt-lr}.

\paragraph{3. Proof of the effective reuse rate scaling law.} With the one-pass and multi-epoch SGD training learning rate set to the near-optimal learning rate obtained above, the proof of \Cref{thm:convex-E_K} proceeds to characterise the behaviour of $E(K,N)$ as $K$ and $N$ vary, yielding the desired scaling relation.
Together, these three components establish \Cref{thm:convex-E_K} and provide a comprehensive description of how reuse efficiency depends on the interplay between $K$ and $N$.